\documentclass[11pt]{article}
\usepackage{amsmath,amssymb,mathtools}
\usepackage[margin=1in]{geometry}
\usepackage{hyperref}
\usepackage{bm}
\usepackage{booktabs}
\usepackage{array}
\usepackage{enumitem}
\usepackage{graphicx}

\setlist{noitemsep,topsep=2pt,parsep=0pt,partopsep=0pt}
\newcolumntype{L}[1]{>{\raggedright\arraybackslash}p{#1}}
\graphicspath{{docs/images/examples/}{../images/examples/}}

\title{Multi-Solver Coupling Strategies \\
       \large WIP design note}
\author{Newton coupling notes}
\date{2026-04-23}

\newcommand{\va}{\bm v_a}
\newcommand{\vb}{\bm v_b}
\newcommand{\uu}{\bm u}
\newcommand{\ubar}{\bar{\bm u}}
\newcommand{\fa}{\bm f_a}
\newcommand{\fb}{\bm f_b}
\newcommand{\lam}{\bm\lambda}
\newcommand{\Ja}{J_a}
\newcommand{\Jb}{J_b}
\newcommand{\Ma}{M_a}
\newcommand{\Mb}{M_b}
\newcommand{\Mhat}{\widehat M_a}
\newcommand{\Dop}{D}
\newcommand{\Dhat}{\widehat D}
\newcommand{\Keff}{K_{\mathrm{eff}}}
\newcommand{\Pmat}{P}
\newcommand{\res}{\bm r}
\newcommand{\normal}{\bm n}

\begin{document}
\maketitle

\section{Scope}

This is a work-in-progress design note. It uses a common notation to describe several
candidate coupling schemes for running two solvers on one simulation step. The text is
not a statement of shipped behavior. In particular, solver composition, model views,
proxy bodies, and ADMM coupling are treated here as proposed design surfaces.

The variants considered below are:
\begin{itemize}
\item collider exchange,
\item one-way coupling,
\item proxy bodies / virtual inertia,
\item linearized ADMM,
\item implicit-interface ADMM.
\end{itemize}
They all start from the same velocity-level interface problem. The main differences
are where the interface law is solved, how action-reaction consistency is obtained, and
what API each scheme requires from the sub-solvers.

Notation map from earlier notes: $(M_1,M_2,H_1,H_2)$ are written here as
$(\Ma,\Mb,\Ja,\Jb)$ for solvers $a$ and $b$. The Delassus operator is denoted by
$\Dop$. The ADMM contact-space preconditioner is denoted by $\Pmat$.

\section{Common interface problem}

Let $\va\in\mathbb R^{n_a}$ and $\vb\in\mathbb R^{n_b}$ be the generalized velocities
owned by solvers $a$ and $b$ during one time step. The matrices $\Ma$ and $\Mb$ denote
the step inertia or the local quadratic model used by each sub-solver. The vectors
$\fa,\fb$ collect explicit forces, controls, and internal terms after the local
time-integration linearization.

The interface relative velocity is
\begin{equation}
    \uu = \Ja\va + \Jb\vb,
    \label{eq:relative-velocity}
\end{equation}
where $\Ja$ and $\Jb$ are contact or attachment Jacobians. They may represent normal
contact velocities, tangential slip velocities, rigid attachment velocities, or other
low-dimensional interface quantities. Bias velocities and stabilization terms can be
folded into a target $\ubar$.

For a hard bilateral velocity constraint $\uu=\ubar$, the monolithic coupled solve is
\begin{equation}
\begin{bmatrix}
\Ma & 0 & -\Ja^\top \\
0 & \Mb & -\Jb^\top \\
-\Ja & -\Jb & 0
\end{bmatrix}
\begin{bmatrix}
\va \\ \vb \\ \lam
\end{bmatrix}
=
\begin{bmatrix}
\fa \\ \fb \\ -\ubar
\end{bmatrix}.
\label{eq:global-kkt}
\end{equation}
The multiplier $\lam$ is the interface impulse in the sign convention
\[
    \Ma\va = \fa + \Ja^\top\lam,\qquad
    \Mb\vb = \fb + \Jb^\top\lam .
\]
For unilateral contact, the third row is replaced by a contact law. In the frictionless
normal direction this is the Signorini condition
\[
    0 \leq u_n \perp \lambda_n \geq 0 .
\]
Friction modifies the local contact law but not the ownership split.

\paragraph{Delassus form.}
Eliminating $\va$ and $\vb$ gives the contact-space relation
\begin{equation}
    \uu = \Dop\lam + \uu^0,\qquad
    \Dop =
        \Ja\Ma^{-1}\Ja^\top
      + \Jb\Mb^{-1}\Jb^\top,\qquad
    \uu^0 =
        \Ja\Ma^{-1}\fa
      + \Jb\Mb^{-1}\fb .
    \label{eq:delassus}
\end{equation}
$\Dop$ is the Delassus operator, or contact-space inverse inertia. Its inverse
$\Keff=\Dop^{-1}$ is the effective mass when it exists.

\paragraph{Finite stiffness form.}
For a finite bilateral stiffness $\kappa$ with energy
$E_c(\uu)=\frac{\kappa}{2}\|\uu-\ubar\|^2$, the monolithic velocity solve can also be
written as
\begin{equation}
\begin{bmatrix}
\Ma + \kappa\Ja^\top\Ja & \kappa\Ja^\top\Jb \\
\kappa\Jb^\top\Ja & \Mb + \kappa\Jb^\top\Jb
\end{bmatrix}
\begin{bmatrix}
\va \\ \vb
\end{bmatrix}
=
\begin{bmatrix}
\fa + \kappa\Ja^\top\ubar \\
\fb + \kappa\Jb^\top\ubar
\end{bmatrix}.
\label{eq:stiffness}
\end{equation}
All variants below avoid directly assembling this full coupled system.

\section{Coupling variants}

\subsection{Collider exchange}

In collider exchange, each sub-solver sees the other side as kinematic collision
geometry. The two sub-solvers advance independently, each with its own contact solve.
This is the least intrusive scheme: it requires only pose and velocity updates for the
foreign collision geometry.

The approximation is that the two sides do not solve one shared interface impulse.
Consequently, the impulses computed by the two contact solvers need not be equal and
opposite. The combined interface complementarity problem is also not enforced as one
problem, and either side can receive kinematic boundary velocities that make its local
contact problem inconsistent. The method is mainly justified when the time step is
small or when the foreign side is intentionally treated as prescribed scenery.

\subsection{One-way coupling}

One-way coupling is the limiting case where side $a$ is prescribed or much heavier than
side $b$. Setting $\Ma^{-1}$ to zero in \eqref{eq:delassus}, side $a$ first advances to
\[
    \va^\star=\Ma^{-1}\fa ,
\]
and side $b$ solves
\begin{equation}
    \uu =
      \Ja\va^\star
    + \Jb\Mb^{-1}\fb
    + \Jb\Mb^{-1}\Jb^\top\lam .
    \label{eq:one-way}
\end{equation}
The interface impulse is scaled by side $b$'s effective inertia. If
$\Mb \ll \Ma$, applying $\Ja^\top\lam$ back to side $a$ produces a negligible change,
so the approximation is consistent with the mass ratio. If the mass ratio assumption is
not valid, the scheme becomes a modeling choice rather than a controlled
approximation.

\subsection{Proxy bodies / virtual inertia}

In the proxy-body variant, objects owned by side $a$ are represented in side $b$ by
replica degrees of freedom with an approximate inertia $\Mhat$. The proxy inertia is
chosen to be easy for side $b$ to solve, for example independent rigid-body blocks or a
diagonal approximation. A scalar relaxation factor can be used to make the proxies
lighter or heavier in the destination solve.

Suppose an impulse $J_{a,k}^\top\lam^k$ was applied to side $a$ in the previous outer
iteration or previous time step, producing
\begin{equation}
    \va^k = \Ma^{-1}\left(\fa + J_{a,k}^\top\lam^k\right).
    \label{eq:lagged-va}
\end{equation}
The proxy solve uses the virtual-inertia approximation
\begin{equation}
    \va
    \approx
    \va^k
    + \Mhat^{-1}\left(\Ja^\top\lam - J_{a,k}^\top\lam^k\right),
    \label{eq:virtual-inertia-update}
\end{equation}
or equivalently
\begin{equation}
    \Mhat(\va-\va^k)=\Ja^\top\lam-J_{a,k}^\top\lam^k.
    \label{eq:virtual-inertia-increment}
\end{equation}
The destination solver then solves
\begin{equation}
\begin{bmatrix}
\Mhat & 0 & -\Ja^\top \\
0 & \Mb & -\Jb^\top \\
-\Ja & -\Jb & 0
\end{bmatrix}
\begin{bmatrix}
\va \\ \vb \\ \lam
\end{bmatrix}
=
\begin{bmatrix}
\Mhat\va^k - J_{a,k}^\top\lam^k \\
\fb \\
-\ubar
\end{bmatrix}.
\label{eq:proxy-kkt}
\end{equation}
The subtraction term $-J_{a,k}^\top\lam^k$ prevents the previously applied impulse from
being counted again. The resulting impulse can be applied to side $a$ in the next time
step, or in an outer fixed-point iteration.

For frozen Jacobians and bilateral constraints, define the proxy Delassus
\[
    \Dhat = \Ja\Mhat^{-1}\Ja^\top + \Jb\Mb^{-1}\Jb^\top .
\]
The fixed-point error propagation for the impulse contains the matrix
\begin{equation}
    \Dhat^{-1}\Ja\left(\Mhat^{-1}-\Ma^{-1}\right)\Ja^\top .
    \label{eq:proxy-convergence}
\end{equation}
A spectral radius below one gives a local contraction. This condition is favored when
$\Mhat$ is close to $\Ma$ in the interface directions, when side $b$ is light relative
to side $a$, or when the proxy inertia is relaxed enough to damp the update.

For a scalar interface with $\Ja=1$, write the virtual-inertia scale as
$s=\Mhat/\Ma$ and the mass ratio as $q=\Ma/\Mb$. Then
\begin{equation}
    r(s,q)
    =
    \frac{\Mhat^{-1}-\Ma^{-1}}{\Mhat^{-1}+\Mb^{-1}}
    =
    \frac{1/s-1}{1/s+q}
    =
    \frac{1-s}{1+qs}.
    \label{eq:proxy-scalar-factor}
\end{equation}
Thus $s=1$ gives exact inertia matching and $r=0$. For $0<s<1$ the scalar map is
contractive for any positive $q$, but the limit $s\to0$ approaches $r\to1$ from
below, so very light proxies are only weakly contractive. For $s>1$ the fixed-point
factor changes sign, and it becomes non-contractive when the overscaled proxy is large
enough relative to the side-$b$ inertia.

The sign and magnitude of Eq.~\eqref{eq:proxy-scalar-factor} are useful as a
two-parameter diagnostic over $s$ and $q$.  The repository keeps the executable
visualization in \texttt{example\_proxy\_coupling\_convergence.py}; the generated
plot is intentionally not checked into the documentation image gallery while
the coupled examples remain experimental.

When the true source-side effective inertia is not available, the proxy mass scale
should be treated as a relaxation parameter rather than as a precise physical mass.
A practical policy is:
\begin{enumerate}
\item build the cheapest available estimate $M_{\mathrm{est}}$ of source inertia at
      the interface: particle mass for particles, directional rigid effective mass
      $1/(\Ja\Ma^{-1}\Ja^\top)$ when available, translational body mass otherwise,
      or a local particle-stencil mass for cloth and soft bodies;
\item initialize $\Mhat=sM_{\mathrm{est}}$ with $s\in[0.5,1]$, using the lower end
      when the estimate is crude and the upper end when the estimate is trusted;
\item measure feedback convergence from harvested forces or impulses, for example
      $e_k=\|\lam_k-\lam_{k-1}\|$ or, in force form,
      $e_k=\|\mathbf f_k-\mathbf f_{k-1}\|$;
\item adapt the scale from the observed residual ratio
      $\hat r=e_k/\max(e_{k-1},\epsilon)$: reduce $s$ when $\hat r$ is large or the
      feedback oscillates, keep it when $\hat r$ is moderate, and increase $s$
      toward one when $\hat r$ is small;
\item clamp the result, for example $s\in[0.05,1]$, so the proxy stays conservative
      and avoids overscaled-inertia sign flips unless a solver-specific effective
      mass estimate justifies $s>1$.
\end{enumerate}
This requires no new mandatory solver interface if adaptation happens between
top-level steps: the shared coupler can use its existing harvested feedback buffers,
update the stored proxy scale, and rescale the destination \texttt{ModelView} mass
overlay before the next step. A runtime API to mutate proxy mass scales, or a small
\texttt{notify\_model\_changed} path for view-local mass overlays, would make this
implementation cleaner and avoid rebuilding the coupled solver. Per-interface
effective-mass hooks would still be useful for better initialization, but they are an
optimization rather than a prerequisite for adaptive proxy relaxation.

\subsection{Linearized ADMM}

ADMM introduces an explicit interface variable and keeps the ownership of physical
degrees of freedom disjoint. The coupling law is represented by an interface energy or
indicator $E_c(\uu)$:
\begin{itemize}
\item hard bilateral attachment: indicator of $\uu=\ubar$;
\item compliant attachment: $E_c(\uu)=\frac{\kappa}{2}\|\uu-\ubar\|^2$;
\item frictionless unilateral contact: indicator of the admissible normal velocity set;
\item frictional contact: local maximum-dissipation Coulomb velocity solve.
\end{itemize}
The split problem is
\begin{equation}
\min_{\va,\vb,\uu}\;
    E_a(\va) + E_b(\vb) + E_c(\uu)
\quad\text{s.t.}\quad
    \uu = \Ja\va + \Jb\vb .
\label{eq:admm-split}
\end{equation}
$E_a$ and $E_b$ are the sub-solvers' per-step objectives, including their internal
forces and contacts.

Let
\[
    \res(\va,\vb,\uu)=\uu-\Ja\va-\Jb\vb .
\]
Using the physical interface impulse $\lam$ as the dual variable and a symmetric
positive-definite contact-space preconditioner $\Pmat$, the augmented Lagrangian is
\begin{equation}
\mathcal L_\rho(\va,\vb,\uu,\lam)
=
E_a(\va)+E_b(\vb)+E_c(\uu)
+ \langle \lam,\,\res\rangle
+ \frac{\rho}{2}\|\res\|_{\Pmat}^2,
\qquad
\|x\|_{\Pmat}^2 = x^\top\Pmat x.
\label{eq:alag}
\end{equation}
A natural choice is $\Pmat\approx\Keff=\Dop^{-1}$, so the penalty is scaled by
contact-space effective mass. Equivalently, one may write
$\Pmat=W^\top W$ with $W\approx\Keff^{1/2}$.

A direct ADMM velocity update would add the Hessian
$\rho\Ja^\top\Pmat\Ja$ to solver $a$ and $\rho\Jb^\top\Pmat\Jb$ to solver $b$.
Linearized ADMM avoids that solver-internal Hessian by evaluating the penalty gradient
at the present iterate and adding a local proximal term. With
\[
    \res^k = \uu^k-\Ja\va^k-\Jb\vb^k,\qquad
    \bm g^k = \lam^k + \rho\Pmat\res^k ,
\]
the parallel velocity sweep is
\begin{align}
\va^{k+1}
&= \arg\min_{\va}\;
    E_a(\va)
    - \langle \Ja^\top\bm g^k,\,\va\rangle
    + \frac{\alpha_a}{2}\|\va-\va^k\|_{\Ma}^2,
    \label{eq:admm-va}\\
\vb^{k+1}
&= \arg\min_{\vb}\;
    E_b(\vb)
    - \langle \Jb^\top\bm g^k,\,\vb\rangle
    + \frac{\alpha_b}{2}\|\vb-\vb^k\|_{\Mb}^2.
    \label{eq:admm-vb}
\end{align}
A sufficient convexity condition is
\[
    \alpha_a\Ma \succeq \rho\,\Ja^\top\Pmat\Ja,\qquad
    \alpha_b\Mb \succeq \rho\,\Jb^\top\Pmat\Jb .
\]

After the velocity sweep, the interface variable is updated by the small separable
problem
\begin{equation}
\uu^{k+1}
=
\arg\min_{\uu}\;
    E_c(\uu)
    + \langle \lam^k,\,\uu\rangle
    + \frac{\rho}{2}
      \|\uu-\Ja\va^{k+1}-\Jb\vb^{k+1}\|_{\Pmat}^2 ,
\label{eq:admm-u}
\end{equation}
followed by the dual update
\begin{equation}
    \lam^{k+1}
    =
    \lam^k
    + \rho\Pmat
      \left(\uu^{k+1}
      - \Ja\va^{k+1}
      - \Jb\vb^{k+1}\right).
\label{eq:admm-dual}
\end{equation}
For a damped quadratic attachment,
\[
    E_c(\uu)=\frac{\kappa}{2}\|\uu-\ubar\|^2
            +\frac{\eta}{2}\|\uu\|^2,
\]
the local update remains diagonal per row:
\[
    \uu^{k+1}
    =
    \frac{\rho W^2(\Ja\va^{k+1}+\Jb\vb^{k+1}) + \kappa\ubar - W\lam^k}
         {\kappa+\eta+\rho W^2}.
\]
The damping coefficient $\eta$ adds resistance to relative interface velocity
without changing the Baumgarte or spring target $\ubar$.
For model-joint box friction with per-axis limits $\bm\tau$, the same local
problem uses the support function $\sum_i \tau_i |u_i|$ and reduces to a
component-wise soft-threshold of the relative angular velocity in the joint
frame.  This is the box analogue of the contact maximum-dissipation solve.
The interface update is a tiny dense solve for quadratic attachments, a clamp for a
frictionless normal contact, or a local maximum-dissipation Coulomb solve for
frictional contact. With isotropic contact weight $W$ and
\[
    \uu^\star = \Ja\va^{k+1} + \Jb\vb^{k+1} - \frac{\lam^k}{\rho W},
    \qquad
    \bar\uu = \uu^\star - u_{\min}\normal ,
\]
the implemented Coulomb update applies the Daviet local map in the normal/tangent
frame:
\[
\uu^{k+1} =
u_{\min}\normal +
\begin{cases}
\bar\uu, & \bar u_n \geq 0,\\
0, & \|\bar\uu_t\| \leq -\mu \bar u_n,\\
\left(1 + \mu \bar u_n / \|\bar\uu_t\|\right)\bar\uu_t,
& \text{otherwise}.
\end{cases}
\]
The middle branch is sticking. The last branch is sliding: the normal velocity is
at the unilateral bound and the dual update places the tangential contact force on
the Coulomb boundary, opposite the slip direction. This solves the Coulomb
maximum-dissipation law rather than projecting an already-computed force onto a
cone.

\subsection{Implicit-interface ADMM}

The full ADMM variant keeps the quadratic penalty inside each sub-solver instead of
linearizing it. The update for side $a$ contains the term
\[
    \frac{\rho}{2}
    \|\uu^k-\Ja\va-\Jb\vb^k\|_{\Pmat}^2 ,
\]
and side $b$ receives the analogous term. This removes the proximal tuning parameters
$\alpha_a,\alpha_b$ and usually reduces the number of ADMM iterations. The cost is a
stronger sub-solver API: each solver must accept implicit point attachments, implicit
springs, or another representation of $\Ja^\top\Pmat\Ja$ within its own solve.

This variant is closest to the monolithic finite-stiffness system
\eqref{eq:stiffness}, while still allowing the two velocity blocks to be minimized by
their native solvers.

\section{API additions relative to upstream/main}

The upstream/main baseline provides a generic
\[
    \texttt{SolverBase.step(state\_in, state\_out, control, contacts, dt)}
\]
interface, external force arrays on \texttt{State} such as \texttt{body\_f} and
\texttt{particle\_f}, model-change notification via \texttt{notify\_model\_changed},
and model/state attribute metadata. It does not provide a public coupled-solver
orchestration layer, model partition views, proxy-body metadata, interface Jacobian
actions, or a standard solver API for effective mass and implicit interface terms.

\subsection{Common APIs}

Most variants benefit from the following additions:
\begin{itemize}
\item \textbf{Model partitioning.} A way to define solver ownership for bodies,
      particles, shapes, and joints. This should include collision ownership:
      internal pairs owned by one sub-solver, cross-owner pairs retained for interface
      coupling, and static geometry assigned explicitly.
\item \textbf{Model views.} A lightweight view or subset object that presents one
      sub-solver with only the attributes it owns, plus controlled overrides such as
      mass, inertia, collision masks, or kinematic state. This avoids cloning whole
      models for every variant.
\item \textbf{State distribution and reconciliation.} Utilities to scatter a global
      state into per-solver states and gather the owned outputs back. The API must
      define which side owns overlapping quantities and how joint, body, and particle
      arrays are merged.
\item \textbf{External-force contract.} A solver-level guarantee that
      \texttt{body\_f} and \texttt{particle\_f} are interpreted as additive external
      loads for that step, and a clear rule for whether solvers read, modify, or
      preserve those arrays.
\item \textbf{Coupling constraint data.} Internal row records derived from model
      joints or collision detection: endpoint type and index, local anchor,
      normal/tangent basis, target velocity or bias, compliance, friction, and
      owner side.
\item \textbf{Jacobian actions.} Kernels or methods for evaluating $Jv$ and applying
      $J^\top f$ for rigid bodies, particles, and mixed endpoint pairs.
      Materializing $J$ should not be required.
\item \textbf{Effective-mass estimates.} Optional methods to evaluate diagonal or
      block-diagonal approximations of $J M^{-1} J^\top$ for contact-space scaling.
\end{itemize}

\subsection{APIs for collider exchange and one-way coupling}

These schemes need the smallest API surface:
\begin{itemize}
\item update kinematic collision geometry from another solver's state;
\item filter collision pairs so that each sub-solver receives the intended collider
      set;
\item optionally accumulate or report reaction forces if one-way feedback is desired.
\end{itemize}
They do not require a shared interface dual variable or repeated in-step iterations.

\subsection{APIs for proxy / virtual-inertia coupling}

Proxy coupling requires additional support:
\begin{itemize}
\item create or expose replica bodies or particles in the destination model without
      treating them as independently owned physical outputs;
\item mark proxy bodies and particles as coupling metadata rather than ordinary
      dynamic or kinematic entities;
\item override proxy mass and inertia in a solver view, including block rigid inertia
      or diagonal approximations;
\item sync proxy poses and velocities from the source side before the destination
      solve;
\item remove previously applied impulses, gravity, or other loads that would otherwise
      be counted twice on proxy degrees of freedom;
\item report contact impulses or accumulated wrenches on proxies after the destination
      solve so that the source side can receive the reaction;
\item optionally run an outer fixed-point loop with convergence metrics on interface
      velocity and impulse.
\end{itemize}

\subsection{APIs for linearized ADMM}

Linearized ADMM can avoid arbitrary coupling Hessians, but it needs iteration-level
control:
\begin{itemize}
\item repeat sub-solver steps from the same begin-of-step state while changing only
      external coupling forces and proximal targets;
\item apply $J^\top\bm g$ as body and particle forces before each ADMM velocity sweep;
\item express the proximal term
      $\frac{\alpha}{2}\|v-v^k\|_M^2$ either as a first-class solver option or through
      mass scaling plus a rewritten pre-step velocity;
\item keep warm-started interface variables $\uu$ and $\lam$ across time steps;
\item expose residuals, stopping criteria, and graph-capture-safe iteration buffers.
\end{itemize}

\subsection{APIs for implicit-interface ADMM}

The implicit variant needs the strongest solver interface:
\begin{itemize}
\item add per-interface implicit springs or attachments to a solver step;
\item represent local Hessian contributions such as $\rho J^\top P J$ without forcing
      callers to assemble solver-internal matrices;
\item expose enough effective-mass information to tune $\rho$ and $P$ robustly;
\item preserve the same interface data model as linearized ADMM so the two ADMM
      variants can share contact detection, Jacobian actions, and dual updates.
\end{itemize}

\section{Implemented public Python API}

The implemented public API lives in \texttt{newton.solvers}; examples and user code
import those public symbols. The implementation modules remain under
\texttt{newton.\_src}.

The implemented split uses a shared coupled-solver base plus algorithm-specific
derived couplers. The shared coupler owns state partitioning, \texttt{ModelView}
construction, state reconciliation, per-entry substeps, and generic coupling hook
helpers. Lagged-impulse proxy coupling lives in \texttt{SolverProxyCoupled};
linearized ADMM lives in \texttt{SolverAdmmCoupled}.

The currently exported coupling symbols are:
\begin{itemize}
\item \texttt{ModelView} and \texttt{SolverCoupled};
\item \texttt{SolverProxyCoupled} and \texttt{SolverAdmmCoupled}.
\end{itemize}

\subsection{Shared coupled-solver API}

\begin{verbatim}
from collections.abc import Callable, Sequence
from dataclasses import dataclass, field
from enum import IntEnum

import newton as nt
import warp as wp
from newton.solvers import SolverBase


class SolverCoupled(SolverBase):
    class ContactPolicy(IntEnum):
        INTERNAL_ONLY = 0
        INTERNAL_AND_STATIC = 1
        NONE = 2

    @dataclass(frozen=True)
    class Entry:
        name: str
        solver: type[SolverBase] | Callable
        bodies: Sequence[int] = ()
        particles: Sequence[int] = ()
        joints: Sequence[int] = ()
        shapes: Sequence[int] = ()
        solver_kwargs: dict = field(default_factory=dict)
        configure_view: Callable[[ModelView], None] | None = None
        contact_policy: int | None = None
        substeps: int = 1
        in_place: bool = False

    def __init__(
        self,
        model: nt.Model,
        entries: Sequence["SolverCoupled.Entry"],
        coupling: object | None = None,
    ) -> None: ...

    def get_solver(self, name: str) -> SolverBase: ...

    def solver(self, name: str) -> SolverBase: ...

    def get_view(self, name: str) -> ModelView: ...


class SolverProxyCoupled(SolverCoupled):
    class ProxyMode(IntEnum):
        LAGGED = 0
        STAGGERED = 1

    @dataclass(frozen=True)
    class Proxy:
        source: str
        destination: str
        bodies: Sequence[int] = ()
        proxy_bodies: Sequence[int] | None = None
        mass_scale: float = 1.0
        mode: int | str = 0
        particles: Sequence[int] = ()
        proxy_particles: Sequence[int] | None = None
        collision_pipeline: Callable[[ModelView], object | None] | None = None
        collide_interval: int | None = None

    @dataclass(frozen=True)
    class CouplingProxy:
        proxies: Sequence["SolverProxyCoupled.Proxy"]
        iterations: int = 1

    def __init__(
        self,
        model: nt.Model,
        entries: Sequence["SolverCoupled.Entry"],
        coupling: "SolverProxyCoupled.CouplingProxy",
    ) -> None: ...

    def get_proxy_contacts(
        self,
        source: str,
        destination: str,
    ) -> nt.Contacts | None: ...
\end{verbatim}

\texttt{SolverCoupled.Entry} describes ownership, sub-solver construction, and
optional entry-local time substeps. The solver callable is constructed as
\texttt{solver(model=view, **solver\_kwargs)}, so it must accept a
\texttt{model} keyword. \texttt{substeps} splits that entry's solve into equal
sub-intervals. \texttt{in\_place=True} aliases the entry's input and output
states and calls \texttt{solver.step(state\_0, state\_0, ...)}; use it only for
solvers that explicitly support same-object input/output stepping. In-place
entries currently require \texttt{substeps=1}, and lagged proxy coupling rejects
in-place source entries because the source begin pose must remain available.
\texttt{SolverCoupled} builds one
\texttt{ModelView} per entry, zeros non-owned body or particle inverse masses in that
view, and reconciles only owned body, particle, and joint
outputs back into the caller's \texttt{state\_out}. The constructor only accepts
explicit \texttt{entries}; older map-based prototype construction paths have been
removed. Derived couplers may keep non-owned endpoints dynamic in a view and mark
them as proxies before sub-solver construction. Proxy coupling is enabled by
constructing \texttt{SolverProxyCoupled} with a \texttt{CouplingProxy} object.

\subsection{\texttt{ModelView}}

\begin{verbatim}
class ModelView:
    def __init__(self, parent: nt.Model, name: str) -> None: ...

    @property
    def name(self) -> str: ...

    @property
    def parent(self) -> nt.Model: ...

    @property
    def overrides(self) -> dict[str, object]: ...

    def state(self, requires_grad: bool | None = None) -> nt.State: ...

    def zero_body_mass(self, body_indices: wp.array) -> None: ...

    def scale_body_mass(self, body_indices: wp.array, factor: float) -> None: ...

    def mark_proxy_bodies(self, body_indices: wp.array) -> None: ...

    def mark_proxy_particles(self, particle_indices: wp.array) -> None: ...

    def zero_particle_mass(self, particle_indices: wp.array) -> None: ...

    def scale_particle_mass(
        self,
        factor: float,
        particle_indices: wp.array | None = None,
    ) -> None: ...
\end{verbatim}

\texttt{ModelView} has copy-on-write overlay semantics for its explicit mutator
methods. Reads fall through to the parent model until a mutator clones the affected
array into the view's overrides. The parent model is not mutated by
\texttt{zero\_body\_mass}, \texttt{scale\_body\_mass},
\texttt{mark\_proxy\_bodies}, \texttt{mark\_proxy\_particles},
\texttt{zero\_particle\_mass}, or \texttt{scale\_particle\_mass}. Direct writes
through a returned Warp array are not intercepted, so callers should use the mutator
methods for view-local edits.

\subsection{Proxy / virtual-inertia coupling}

\begin{verbatim}
solver = SolverProxyCoupled(
    model,
    entries=[rigid_entry, soft_entry],
    coupling=SolverProxyCoupled.CouplingProxy(
        proxies=[
            SolverProxyCoupled.Proxy(
                source="rigid",
                destination="soft",
                bodies=rigid_body_ids,
                proxy_bodies=rigid_proxy_body_ids,
                mass_scale=0.25,
                mode=SolverProxyCoupled.ProxyMode.LAGGED,
                collision_pipeline=lambda model: CollisionPipeline(model),
                collide_interval=1,
            )
        ]
    ),
)
\end{verbatim}

\texttt{proxy\_bodies=None} maps each source body to the same body index in the
destination view; otherwise \texttt{proxy\_bodies} must have the same length as
\texttt{bodies}. \texttt{proxy\_particles=None} does the same for source particles.
\texttt{mass\_scale} scales proxy body mass/inertia and proxy particle mass in the
destination view. \texttt{mode} accepts either the \texttt{ProxyMode} enum value or
the strings \texttt{"lagged"} and \texttt{"staggered"}. \texttt{LAGGED} syncs the
source begin-of-step state and end velocity, then rewinds destination proxy
velocities by the previously applied feedback wrench or particle force, public
force input, and gravity unless the destination solver provides a custom body or
particle proxy velocity-rewind hook.
\texttt{STAGGERED} syncs the source end state and end velocity and skips the
generic rewind. If \texttt{collision\_pipeline} is supplied, the proxy coupler
calls the factory with the destination \texttt{ModelView}. If the factory returns
\texttt{None}, no proxy-local collision pipeline is installed and the destination
solve receives the outer-level contacts supplied to \texttt{step()}. Otherwise
the coupler allocates a persistent contacts buffer from the returned pipeline and
refreshes those contacts before the destination solve every
\texttt{collide\_interval} proxy passes. A \texttt{None} interval means every
proxy pass when a custom pipeline is supplied; with no custom pipeline, the
coupler keeps the previous behavior and uses the contacts supplied to
\texttt{step()}. \texttt{get\_proxy\_contacts()} exposes the internal contact
buffer for diagnostics and rendering; it returns \texttt{None} when no
proxy-local pipeline is installed.
Destination solvers may also mark
\texttt{BODY\_PROXY\_HARVEST} or \texttt{PARTICLE\_PROXY\_HARVEST} as custom to
override generic momentum harvesting; otherwise the shared coupler estimates
feedback from proxy momentum change. Custom harvest hooks receive the prepared
destination input state, the destination output state, and the contacts used for
the proxy solve, so contact-native solvers can report feedback from final contact
forces rather than from aggregate solver forces or arbitrary momentum changes.
VBD uses this path to harvest rigid-rigid contact forces through its per-contact
collector and body-particle contact forces through a matching soft-contact
reducer. MPM has two particle-proxy paths. Pure proxy particles stay active in
the MPM transfer flags, so they participate in the normal P2G/G2P momentum
transfer, but are masked out of material volume, stress, strain, and constitutive
updates. Their feedback is harvested from the transfer-induced momentum change
with gravity removed from the lagged velocity rewind and from the final force.
Deformable collider particles registered through \texttt{collider\_particle\_ids}
remain inactive in both transfer and material terms and use the private collider
impulse collector instead. The \texttt{iterations} field repeats this proxy
relaxation pass over the same top-level time interval.
The implemented generic proxy loop currently supports at most two solver entries;
larger proxy-coupled configurations should use a more explicit scheduling layer.
Within that two-entry limit, the loop groups body and particle mappings by
\texttt{(source, destination)} and performs one source step and one destination
step per solver pair per proxy iteration. A single \texttt{Proxy} may therefore
carry both \texttt{bodies} and \texttt{particles}; both mappings are synchronized,
velocity-rewound, and harvested around the same destination solve.

Example usage:
\begin{verbatim}
rigid = SolverProxyCoupled.Entry(
    name="rigid",
    solver=SolverMuJoCo,
    bodies=robot_bodies,
    shapes=robot_shapes,
    joints=robot_joints,
    substeps=4,
)
soft = SolverProxyCoupled.Entry(
    name="soft",
    solver=SolverVBD,
    bodies=cloth_bodies,
    particles=cloth_particles,
    shapes=cloth_shapes,
)

solver = SolverProxyCoupled(
    model,
    entries=[rigid, soft],
    coupling=SolverProxyCoupled.CouplingProxy(
        proxies=[
            SolverProxyCoupled.Proxy(
                source="rigid",
                destination="soft",
                bodies=robot_bodies,
                proxy_bodies=robot_proxy_bodies,
                mass_scale=0.25,
                mode=SolverProxyCoupled.ProxyMode.LAGGED,
            )
        ],
    ),
)
\end{verbatim}


\subsection{Linearized ADMM coupling}

\begin{verbatim}
class SolverAdmmCoupled(SolverCoupled):
    @classmethod
    def register_custom_attributes(cls, builder: nt.ModelBuilder) -> None: ...

    @classmethod
    def add_body_particle_attachment(
        cls,
        builder: nt.ModelBuilder,
        body: int,
        particle: int,
        *,
        body_point: tuple[float, float, float] | wp.vec3 = (0.0, 0.0, 0.0),
        stiffness: float = 1.0e4,
        damping: float = 0.0,
        enabled: bool = True,
    ) -> int: ...

    @dataclass(frozen=True)
    class CouplingAdmm:
        iterations: int = 5
        rho: float = 1.0
        gamma: float = 0.0
        baumgarte: float = 0.0
        joint_stiffness: float = 1.0e4
        joint_damping: float = 0.0
        joint_angular_stiffness: float = 1.0e4
        joint_angular_damping: float = 0.0
        contact_detection: bool = False
        contact_distance: float | None = None
        contact_friction: float = 0.0
        contact_detection_margin: float = 0.01
        particle_contact_detection_margin: float | None = None

    def __init__(
        self,
        model: nt.Model,
        entries: Sequence[SolverCoupled.Entry],
        coupling: "SolverAdmmCoupled.CouplingAdmm",
    ) -> None: ...
\end{verbatim}

\texttt{SolverAdmmCoupled} no longer accepts explicit endpoint-interface
records. The ADMM constraints are discovered from the shared model:
\begin{itemize}
    \item cross-solver model joints become quadratic ADMM attachment rows;
    \item custom \texttt{coupling:body\_particle\_attachment} rows become
    quadratic rigid-body to particle attachment rows;
    \item particle-shape, rigid-rigid, and particle-particle contacts come from
    model-derived collision candidates when \texttt{contact\_detection=True}.
\end{itemize}

For cross-solver joints, the two connected bodies must be owned by different
\texttt{SolverCoupled.Entry} objects, and the joint itself must be left unowned
by sub-solvers so it is not solved twice. The current generic ADMM joint path
maps \texttt{JointType.BALL} to translational anchor-coincidence rows and
\texttt{JointType.FIXED} to translational plus angular rows. Other cross-solver
joint types are rejected until their constrained subspace is represented
exactly in the ADMM local solve. The scalar \texttt{joint\_stiffness},
\texttt{joint\_damping}, \texttt{joint\_angular\_stiffness}, and
\texttt{joint\_angular\_damping} parameters are the quadratic local-solve
stiffnesses and damping coefficients for those model-derived rows.
For \texttt{JointType.BALL}, nonzero \texttt{model.joint\_friction} entries
create an additional angular dry-friction row.  The row measures relative
angular velocity in the child joint frame and applies a per-axis box
maximum-dissipation solve, matching the model's three stored ball-joint
friction DOFs.

\texttt{SolverAdmmCoupled.add\_body\_particle\_attachment} is the convenience
path for deformable-to-rigid attachments that cannot be represented by a model
joint because one endpoint is a particle. It registers a model custom frequency
\texttt{coupling:body\_particle\_attachment} with SoA attributes for the body
index, particle index, body-local point, stiffness, damping, and enabled flag, then
appends one row. During construction, rows whose body and particle endpoints are
owned by different entries become quadratic ADMM attachment rows. Rows with
unowned or same-entry endpoints are ignored, matching the model-joint path's
"only cross-solver constraints are owned by the coupler" rule. The helper is
equivalent to calling \texttt{register\_custom\_attributes} and filling the
custom values directly, so USD or importer paths can author the same attributes
without using the Python helper.

\texttt{contact\_detection=True} enables model-derived ADMM contacts. The
coupler owns a private collision pipeline for shape-based contact detection and
builds contact groups from solver-entry ownership: particle-shape rows are
generated between particle entries and shapes on bodies owned by another entry
whose shape flags include particle collision; rigid-rigid rows are generated
from cross-entry shape contact pairs; particle-particle rows are generated for
cross-entry particle sets through a private hash-grid pass. When
\texttt{contact\_distance} is \texttt{None}, rigid contacts use collision
margins, particle-shape contacts use particle radii, and particle-particle
contacts use radius sums. Otherwise \texttt{contact\_distance} is the admissible
signed gap for generated rows. \texttt{contact\_friction} sets the isotropic
Coulomb coefficient. The hash-grid contact stream is an internal fixed-capacity
record used only to refresh persistent ADMM rows and report normal force/impulse
inside the coupler. Contact paths keep fixed-capacity device arrays and
warm-start \texttt{u} / \texttt{lambda} by stable contact keys when contacts
persist from one step to the next.

A possible future extension is to make contact streams the input contract for
proxy momentum harvesting itself: stream rows would encode source/proxy endpoint
pairs, the proxy step would snapshot destination velocities, and the harvester
would write feedback directly through those rows. That design is different from
mirroring the current dense momentum-harvest result into streams after the
harvest has already happened. The private detection contacts are not forwarded
to sub-solvers; callers may still pass their own \texttt{Contacts} buffer for
ordinary solver contacts. \texttt{gamma} controls the proximal term: when it is
positive, \texttt{SolverAdmmCoupled} scales masses in each \texttt{ModelView},
notifies sub-solvers that inertial properties changed, and shifts the
begin-of-iteration velocities toward the previous ADMM iterate.

Model-joint attachment usage:
\begin{verbatim}
link_joint = builder.add_joint_revolute(parent=-1, child=link)
payload_joint = builder.add_joint_ball(
    parent=link,
    child=payload,
    parent_xform=wp.transform(p=wp.vec3(link_hx, 0.0, 0.0)),
    child_xform=wp.transform(p=wp.vec3(-payload_hx, 0.0, 0.0)),
)

rigid = SolverCoupled.Entry(
    name="mjc",
    solver=SolverMuJoCo,
    bodies=[link],
    joints=[link_joint],      # sub-solver-owned world joint
)
soft = SolverCoupled.Entry(
    name="vbd",
    solver=SolverVBD,
    bodies=[payload],         # payload_joint is intentionally unowned
)

solver = SolverAdmmCoupled(
    model,
    entries=[rigid, soft],
    coupling=SolverAdmmCoupled.CouplingAdmm(
        iterations=8,
        rho=1.0,
        gamma=0.1,
        baumgarte=0.2,
        joint_stiffness=1.0e4,
        joint_damping=10.0,
    ),
)
\end{verbatim}

Body-particle attachment usage:
\begin{verbatim}
ball = builder.add_body(...)
center_particle = particle_start + center_j * (dim_x + 1) + center_i

SolverAdmmCoupled.add_body_particle_attachment(
    builder,
    body=ball,
    particle=center_particle,
    body_point=(0.0, 0.0, 0.0),
    stiffness=50.0,
    damping=1.0,
)

solver = SolverAdmmCoupled(
    model,
    entries=[
        SolverCoupled.Entry(name="mjc", solver=SolverMuJoCo, bodies=[ball]),
        SolverCoupled.Entry(
            name="vbd",
            solver=SolverVBD,
            particles=list(range(model.particle_count)),
        ),
    ],
    coupling=SolverAdmmCoupled.CouplingAdmm(iterations=8, rho=50.0),
)
\end{verbatim}

Collision-detected particle-shape contact usage:
\begin{verbatim}
solver = SolverAdmmCoupled(
    model,
    entries=[drop, tray],
    coupling=SolverAdmmCoupled.CouplingAdmm(
        iterations=12,
        rho=45.0,
        baumgarte=0.1,
        contact_detection=True,
        contact_distance=0.04,
        contact_detection_margin=0.08,
    ),
)
\end{verbatim}

Collision-detected rigid-rigid contact usage:
\begin{verbatim}
solver = SolverAdmmCoupled(
    model,
    entries=[box_entry, plane_entry],
    coupling=SolverAdmmCoupled.CouplingAdmm(
        iterations=30,
        rho=5.0,
        gamma=0.2,
        baumgarte=0.03,
        contact_detection=True,
        contact_friction=0.45,
    ),
)
\end{verbatim}

Collision-detected particle-particle contact usage:
\begin{verbatim}
solver = SolverAdmmCoupled(
    model,
    entries=[cloth_a, cloth_b],
    coupling=SolverAdmmCoupled.CouplingAdmm(
        iterations=10,
        rho=30.0,
        baumgarte=0.2,
        contact_detection=True,
        contact_detection_margin=0.02,
    ),
)
\end{verbatim}

\subsection{No single-pair convenience wrappers}

The implementation intentionally avoids exported MuJoCo+VBD, MuJoCo+MPM, or
MuJoCo+VBD-ADMM wrapper classes. The examples construct
\texttt{SolverProxyCoupled}, \texttt{SolverAdmmCoupled}, or the base
\texttt{SolverCoupled} directly.  Example-only sharing was also removed: each
example now carries its own rigid-solver selection and graph-capture setup so
the framework behavior is visible at the construction site.  MPM proxy examples
that use deformable colliders rely on \texttt{SolverImplicitMPM}'s default
constructor-time collider setup against the entry \texttt{ModelView}, so proxy
mass relaxation is visible without an extra example-side setup call. The
XPBD-MPM particle example instead exercises transfer-active pure proxy particles,
with no triangles and no collider-particle registration. This keeps the default
coupler path under test for the common two-solver cases and makes any missing
generic hook visible.

\subsection{Optional \texttt{SolverBase} coupling hooks}

The implemented coupling hooks keep \texttt{step()} source-compatible. A solver
reports, per hook, whether the model/state fallback is adequate, a custom
implementation is required, or the requested hook is unsupported. Omitted hooks
default to the fallback path.

\begin{verbatim}
class SolverBase:
    class CouplingEndpointKind(IntEnum):
        BODY = 0
        PARTICLE = 1

    class CouplingInputStateFlags(IntFlag):
        BODY_Q = 1 << 0
        BODY_QD = 1 << 1
        PARTICLE_Q = 1 << 2
        PARTICLE_QD = 1 << 3
        JOINT_Q = 1 << 4
        JOINT_QD = 1 << 5
        BODY_F = 1 << 6
        PARTICLE_F = 1 << 7
        JOINT_F = 1 << 8
        ITERATION_RESTART = 1 << 9

        BODY = BODY_Q | BODY_QD
        PARTICLE = PARTICLE_Q | PARTICLE_QD
        JOINT = JOINT_Q | JOINT_QD
        FORCE = BODY_F | PARTICLE_F | JOINT_F
        ALL = BODY | PARTICLE | JOINT | FORCE

    class CouplingHooks(IntFlag):
        BODY_PROXY_REWIND_VELOCITY = 1 << 0
        PARTICLE_PROXY_REWIND_VELOCITY = 1 << 1
        BODY_PROXY_HARVEST = 1 << 2
        PARTICLE_PROXY_HARVEST = 1 << 3
        EFFECTIVE_MASS_DIAGONAL = 1 << 4
        EFFECTIVE_MASS_BLOCK = 1 << 5
        NOTIFY_INPUT_STATE_UPDATE = 1 << 6
        PROXY_CONTACT_PREPARE = 1 << 7

    class CouplingCapability(Enum):
        DEFAULT = "default"
        CUSTOM = "custom"
        UNSUPPORTED = "unsupported"

    CouplingCapabilities = dict[CouplingHooks, CouplingCapability]

    def coupling_capabilities(
        self,
    ) -> "SolverBase.CouplingCapabilities":
        return {}

    def coupling_capability(
        self,
        hook: "SolverBase.CouplingHooks",
    ) -> "SolverBase.CouplingCapability":
        return self.coupling_capabilities().get(
            hook,
            SolverBase.CouplingCapability.DEFAULT,
        )

    def coupling_eval_effective_mass(
        self,
        endpoint_kind: wp.array[int],
        endpoint_index: wp.array[int],
        endpoint_local_pos: wp.array[wp.vec3],
        out: wp.array[float],
    ) -> None: ...

    def coupling_eval_effective_mass_block(
        self,
        endpoint_kind: wp.array[int],
        endpoint_index: wp.array[int],
        endpoint_local_pos: wp.array[wp.vec3],
        out_mass: wp.array[float],
        out_inertia: wp.array[wp.mat33] | None = None,
    ) -> None: ...

    def coupling_notify_input_state_update(
        self,
        state: nt.State,
        flags: CouplingInputStateFlags | int,
        dt: float = 0.0,
    ) -> None: ...

    def coupling_harvest_proxy_wrenches(
        self,
        body_local_to_proxy_global: wp.array[int],
        out_body_f: wp.array[wp.spatial_vector],
        *,
        state: nt.State | None = None,
        state_out: nt.State | None = None,
        contacts: nt.Contacts | None = None,
        dt: float = 0.0,
    ) -> None: ...

    def coupling_rewind_proxy_body_velocity(
        self,
        body_local_to_proxy_global: wp.array[int],
        state: nt.State,
        coupling_forces: wp.array[wp.spatial_vector],
        dt: float,
    ) -> None: ...

    def coupling_prepare_proxy_contacts(
        self,
        state: nt.State,
        contacts: nt.Contacts | None,
        *,
        contacts_freshly_detected: bool = False,
    ) -> nt.Contacts | None: ...

    def coupling_harvest_proxy_particle_forces(
        self,
        particle_local_to_proxy_global: wp.array[int],
        out_particle_f: wp.array[wp.vec3],
        *,
        state: nt.State | None = None,
        state_out: nt.State | None = None,
        contacts: nt.Contacts | None = None,
        dt: float = 0.0,
    ) -> None: ...

    def coupling_rewind_proxy_particle_velocity(
        self,
        particle_local_to_proxy_global: wp.array[int],
        state: nt.State,
        coupling_forces: wp.array[wp.vec3],
        dt: float,
    ) -> None: ...

\end{verbatim}

The default implementation is the public model/state path. Body and particle
coupling forces are written directly into \texttt{state.body\_f} and
\texttt{state.particle\_f}; force injection is no longer a solver hook.  The maps
used by the couplers are dense local-to-global index maps.  For example, the
default body path loops over local body ids \texttt{b} and, when
\texttt{body\_local\_to\_global[b] >= 0}, adds
\texttt{body\_f[body\_local\_to\_global[b]]} into \texttt{state.body\_f[b]}.
Normal ADMM force buffers use the identity map. Proxy feedback uses a
source-local-to-global-proxy map, so harvested forces stored at the global proxy
id are applied to the source solver's local id without changing the indexing
convention of the feedback buffer.

\texttt{ITERATION\_RESTART} is a contextual bit on
\texttt{coupling\_notify\_input\_state\_update}.  ADMM and proxy coupling set it
when they restart a solver input state for another relaxation pass over the same
top-level time step.  Solvers with private step history can then realign that
history with the public input state without treating the reset as a physical
teleport.  VBD uses this to reset rigid \texttt{body\_q\_prev} on repeated
coupling iterations, while ordinary proxy pose updates still use the proxy
velocity-target path.

Model-view mass and inertia overrides are also not coupling hooks.  Couplers
mutate the destination \texttt{ModelView} directly: body proxy virtual inertia
uses \texttt{set\_body\_inertial\_properties}, and particle proxy virtual mass
uses \texttt{set\_particle\_mass}.  After such post-construction mutations the
coupler calls the ordinary solver \texttt{notify\_model\_changed} path, using
\texttt{BODY\_INERTIAL\_PROPERTIES} for body mass/inertia changes and a model
properties notification for particle mass changes. Solvers with private model
caches should therefore handle these changes in \texttt{notify\_model\_changed}
rather than through a coupling-specific override hook. ADMM uses the same
\texttt{ModelView} mutation mechanism for the proximal inertia term: when
\texttt{gamma>0}, it scales owned body and particle masses by \texttt{1+gamma}
before sub-solver construction and notifies solvers after construction.

\texttt{NOTIFY\_INPUT\_STATE\_UPDATE} is the generic state-side notification
hook. The fallback is no extra action because the coupler has already written
the updated values directly into the public \texttt{State}. Custom solvers use
this hook to synchronize private input buffers or to translate public state
updates into solver-native history. \texttt{SolverProxyCoupled} calls it after
state distribution, source force-buffer updates, source-to-proxy synchronization,
and lagged proxy velocity rewinds. \texttt{SolverAdmmCoupled} calls it at the
beginning of each ADMM iteration after restoring the iteration state, and after
writing the proximal velocity target
\texttt{(v\_n + gamma v\_k)/(1+gamma)} into \texttt{body\_qd},
\texttt{particle\_qd}, and \texttt{joint\_qd}. There is therefore no separate
ADMM-specific proximal-inertia or proximal-velocity hook; proximal inertia uses
model-view mass scaling, and proximal velocity uses this generic state-update
notification hook.

\texttt{BODY\_PROXY\_REWIND\_VELOCITY} and
\texttt{PARTICLE\_PROXY\_REWIND\_VELOCITY} let a destination solver rewind
synchronized proxy velocities before its step. MPM uses these hooks for collider
velocity preparation because its body and deformable-particle colliders are
prescribed rather than integrated from public force buffers. For pure
transfer-active particle proxies, the same MPM hook also subtracts gravity from
the lagged input velocity so gravity is not harvested as coupling feedback.
Without custom capabilities, \texttt{SolverProxyCoupled} applies the generic
lagged velocity rewind from the destination \texttt{ModelView}, subtracting
previously applied feedback, public force inputs, and gravity.
VBD instead uses \texttt{NOTIFY\_INPUT\_STATE\_UPDATE} for smooth proxy-body
teleport handling and the generic rewind for lagged feedback.
\texttt{BODY\_PROXY\_HARVEST} and \texttt{PARTICLE\_PROXY\_HARVEST} let a solver
report feedback from private buffers or impulses; otherwise proxy feedback is
estimated from destination proxy momentum change.
For body and particle proxies, \texttt{state} is the prepared destination input
state, \texttt{state\_out} is the destination output state after the proxy solve,
and \texttt{contacts} is the contact set passed to that solve. This lets custom
harvesters report contact-only feedback at the accepted configuration. VBD uses
this to avoid harvesting the final AVBD body force buffer, which may include
non-contact terms or only the last nonlinear force assembly, and instead reduces
the final rigid-rigid and body-particle contact forces onto proxy bodies.
Proxy rewind and harvest hooks receive a dense local-to-global-proxy map.
The local id is a body or particle index in the destination solver's
\texttt{ModelView}; proxy entries map to the corresponding global proxy id in
the parent model, and non-proxy entries are \texttt{-1}.  State arrays and
solver-private buffers remain local-indexed.  The \texttt{coupling\_forces},
\texttt{out\_body\_f}, and \texttt{out\_particle\_f} buffers are indexed by
global proxy ids.
\texttt{SolverProxyCoupled} maps the proxy-indexed feedback back to source
local ids through the public force-buffer path when injecting it into the
driving solver.
\texttt{EFFECTIVE\_MASS\_DIAGONAL} and \texttt{EFFECTIVE\_MASS\_BLOCK} are
query hooks, not mutation hooks: they let couplers ask a solver for endpoint
effective mass. Endpoint batches are passed as structure-of-arrays Warp
buffers: \texttt{endpoint\_kind}, \texttt{endpoint\_index}, and
\texttt{endpoint\_local\_pos}. The arrays all have the same length as the
output mass buffer. \texttt{endpoint\_kind} stores body/particle kind values,
\texttt{endpoint\_index} stores the local model-view body or particle id, and
\texttt{endpoint\_local\_pos} stores body-frame points for body endpoints
(\texttt{0} for particles). The diagonal hook defaults to model mass/inertia.
The block hook defaults to model mass and full body inertia, or to the
diagonal/model fallback only when the caller can accept a scalarized body
estimate. Proxy coupling uses the block hook for full virtual body inertia and
the diagonal hook for particle proxy mass; ADMM uses diagonal endpoint masses
for interface/contact weighting and falls back to model masses when no custom
query is supplied.

\texttt{coupling\_prepare\_proxy\_contacts} is called for body-proxy destination
solves so the destination contact set can filter proxy contacts outside its
ownership/feedback boundary. Proxy-local collision generation is configured on
\texttt{SolverProxyCoupled.Proxy} rather than through this hook; the hook receives
\texttt{contacts\_freshly\_detected=True} only when that proxy collision pipeline
has just run. VBD uses this to filter newly detected proxy contacts and to call
\texttt{set\_rigid\_history\_update} on the same cadence as contact refreshes.
XPBD additionally treats view-local \texttt{ParticleFlags.PROXY} particles as
contact-only proxies, filtering proxy-proxy and proxy-static particle contacts
while retaining owned-particle versus proxy contact.

\section{Comparison}

\begin{center}
\small
\begin{tabular}{@{}L{0.18\linewidth}L{0.22\linewidth}L{0.19\linewidth}L{0.28\linewidth}@{}}
\toprule
Variant & Interface solve & API cost & Main tradeoff \\
\midrule
Collider exchange
& independent kinematic-contact solves
& low
& easiest to stage, but action-reaction and complementarity are not solved as one
interface problem \\
\addlinespace
One-way
& one side prescribes motion, the other solves contact
& low
& controlled when one side is effectively infinite mass; otherwise intentionally
asymmetric \\
\addlinespace
Proxy / virtual inertia
& destination solver solves contacts against finite-mass replicas
& medium
& reuses a native contact solver, but stability depends on proxy inertia and force
feedback is lagged unless fixed-point iterations are used \\
\addlinespace
Linearized ADMM
& external-force sub-solver sweeps plus local interface proximal step
& medium-high
& symmetric in ownership and gives residuals/dual impulses, but needs in-step
iterations and proximal tuning \\
\addlinespace
Implicit-interface ADMM
& sub-solvers include implicit interface quadratic terms
& high
& fewer ADMM iterations and less proximal tuning, but requires implicit attachment or
Hessian support in each sub-solver \\
\bottomrule
\end{tabular}
\end{center}

The schemes therefore occupy different points on an implementation spectrum. Collider
exchange and one-way coupling minimize new API surface. Proxy coupling adds replica
metadata and force harvesting in exchange for reusing one solver's contact machinery.
Linearized ADMM adds a shared interface iteration without requiring arbitrary Hessians
inside sub-solvers. Implicit-interface ADMM asks for the richest solver API and is the
closest split form to a monolithic coupled solve.

\end{document}
